\section{完整案例：模型对比}
这张表使用两个独立的列范围：准确率最大值为 100，延迟最大值为 50。
条越长只表示数值越大；准确率通常越高越好，延迟通常越低越好。
在下面的宽幅案例中，效果放在上方、源码放在下方，避免缩小表格字体。

\begin{example}[text and listing,sidebyside=false]{示例 68：准确率与延迟}
\begingroup
\gradbarssetup{
  width=27mm, height=1.5ex,
  label width=3.6em, precision=1
}
\renewcommand{\arraystretch}{1.65}
\begin{tabular}{@{}lcc@{}}
  \toprule
\multicolumn{1}{c}{Model} & \multicolumn{1}{c}{Accuracy (\%)} & \multicolumn{1}{c}{Latency (ms)}\\
  \midrule
  Baseline
    & \gradbar{72.4}
    & \gradbar[theme=teal,max=50]{18.2} \\
  Compact
    & \gradbar{81.6}
    & \gradbar[theme=teal,max=50]{12.5} \\
  Transformer
    & \gradbar{89.3}
    & \gradbar[theme=teal,max=50]{42.8} \\
  \textbf{Proposed}
    & \gradbar{92.7}
    & \gradbar[theme=teal,max=50]{24.1} \\
  \bottomrule
\end{tabular}
\endgroup
\end{example}

\subsection{嵌入自己的论文}
在主文档导言区加载 \key{booktabs} 和 \key{gradbars}，然后复制上述代码即可。
需要标题和交叉引用时，再在外层使用普通 \key{table} 环境、
\key{\caption} 和 \key{\label}。
同一列应保持相同的最大值、轨道宽度和标签槽宽。
